\section{Feature Attribution: SHAP and LIME}

\begin{frame}{Baseline: Permutation Importance}

    \begin{block}{Recipe}
        \begin{enumerate}
            \item Measure the model's score on the validation set.
            \item Shuffle one feature column, breaking its relationship with the target.
            \item Re-score. The \textbf{drop} is that feature's importance.
            \item Repeat several times and average; then restore the column and move on.
        \end{enumerate}
    \end{block}

    \vspace{0.7em}

    \begin{columns}[T]
        \begin{column}{0.48\textwidth}
            \textbf{Why start here}
            \begin{itemize}\small
                \setlength\itemsep{0.3em}
                \item Model-agnostic, ten lines of code, needs only \texttt{predict()}.
                \item Measures importance \textbf{for the metric you care about}.
                \item Computed on \emph{held-out} data, so it reflects generalisation, not fit.
            \end{itemize}
        \end{column}
        \begin{column}{0.48\textwidth}
            \textbf{Where it misleads}
            \begin{itemize}\small
                \setlength\itemsep{0.3em}
                \item \textbf{Correlated features split the credit}, and each looks unimportant
                      because the other covers for it.
                \item Shuffling creates impossible rows (age 8, income \$200k), so you evaluate
                      the model off its data manifold.
                \item It is global only --- it says nothing about a single prediction.
            \end{itemize}
        \end{column}
    \end{columns}
\end{frame}

\begin{frame}{A Trap: Built-in Tree ``Feature Importance''}

    \texttt{RandomForest.feature\_importances\_} is \textbf{mean decrease in impurity} (MDI),
    accumulated on the \emph{training} data.

    \vspace{0.8em}

    \begin{itemize}
        \setlength\itemsep{0.5em}
        \item It is biased towards \textbf{high-cardinality} features: a continuous feature or a
              random ID column offers many split points, so it looks informative even when it is
              pure noise.
        \item It is computed on training data, so it rewards features the model used to overfit.
        \item Classic demonstration: add a column of random integers to your dataset and it can
              outrank real predictors in MDI --- while permutation importance on held-out data
              correctly gives it roughly zero.
    \end{itemize}

    \vspace{0.8em}
    \begin{block}{}
        \centering
        \small Rule: never report \texttt{feature\_importances\_} as ``the model's feature
        importance'' without saying which definition you used. Prefer permutation importance
        or SHAP on held-out data.
    \end{block}
\end{frame}

\begin{frame}{SHAP: Attribution as a Fair Payout}

    \textbf{The analogy.} A team of players cooperates and earns a payout. How much of it does each
    player deserve? Here the ``players'' are the features of one row, and the ``payout'' is
    $f(x) - \mathbb{E}[f(X)]$: how far this prediction sits from the average prediction.

    \vspace{0.8em}

    \begin{block}{Shapley value (Shapley, 1953)}
        Consider adding feature $i$ to a coalition $S$ of features already present, and record the
        change in the prediction. The Shapley value $\phi_i$ is the \textbf{average of that change
        over all possible orderings} in which features could be added:
        \[
            \phi_i \;=\; \sum_{S \subseteq N \setminus \{i\}}
            \frac{|S|!\,\bigl(|N|-|S|-1\bigr)!}{|N|!}
            \Bigl[\, v(S \cup \{i\}) - v(S) \,\Bigr]
        \]
    \end{block}

    \vspace{0.5em}

    \begin{itemize}\small
        \setlength\itemsep{0.3em}
        \item The averaging over orderings is what makes it symmetric: no feature is privileged by
              being ``first''.
        \item There are $2^{|N|}$ subsets, so the exact computation is exponential.
              Everything practical in SHAP is an efficient approximation of this quantity.
    \end{itemize}
\end{frame}

